\section{Related Work}
\label{sec:related}

\paragraph{Test-time compute scaling.}
The idea that inference-time compute can substitute for model size was demonstrated
empirically by \citet{brown2024large} via best-of-$N$ sampling and later formalized
by \citet{snell2024scaling}, who showed that for hard reasoning problems, verifier-guided
search is preferable to repeated independent sampling. \citet{wu2024inference} provides a
compute-optimal analysis mapping problem difficulty to optimal search strategy.
\citet{openai2024o1} introduced chain-of-thought reasoning with extended thinking time,
and subsequent work by \citet{guo2025deepseek} and \citet{qwq2024} has extended this
to open-weight models. Our work is orthogonal to the choice of base model or search
strategy; we provide a principled framework for reducing verification overhead within
any such system.

\paragraph{Process reward models and verifiers.}
\citet{cobbe2021training} trained outcome reward models (ORMs) to re-rank sampled
solutions. \citet{lightman2023lets} introduced process reward models (PRMs) that score
intermediate reasoning steps, demonstrating substantially better performance on the
MATH benchmark \citep{hendrycks2021math}. \citet{setlur2024rewarding} proved
optimality of PRMs under a fixed compute budget and introduced advantage-based PRMs.
\citet{wang2024math} showed that self-consistency and process rewards can be combined.
These works motivate PRMs as the verifier of choice; DVA-MCTS is compatible with any
verifier satisfying our Lipschitz assumption.

\paragraph{Monte Carlo Tree Search for LLMs.}
MCTS was adapted for LLM reasoning by \citet{yao2023tree} (Tree of Thoughts) and
further developed in \citet{chen2024alphamath,zhang2024llamaberry,liu2023don}.
\citet{guo2025deepseek} integrates MCTS with PRMs in a reinforcement learning
framework. Unlike these works, we do not modify the search tree structure; instead,
we provide an adaptive calling policy for the verifier oracle that reduces its cost.

\paragraph{Adaptive computation.}
Reducing computation by skipping expensive operations in learned models has a rich
history: early-exit classifiers \citep{teerapittayanon2016branchynet}, adaptive depth
transformers \citep{graves2016adaptive,elbayad2020depth}, and speculative decoding
\citep{leviathan2023fast} all reduce cost while preserving output quality.
Most related is \citet{liu2024speculative}, who applies speculative execution ideas to
reasoning chains. Our work differs by providing a \emph{regret guarantee} for the
online decision to invoke or skip the verifier, rather than an empirical speedup.

\paragraph{Online resource allocation and bandits.}
Dynamic verifier allocation is formally related to online learning under a compute
budget constraint. The UCB algorithm of \citet{auer2002finite} achieves
$O(\sqrt{T \log K})$ regret in stochastic multi-armed bandits; our setting generalizes
this to a tree-structured action space with varying arm costs. Budget-limited bandit
problems have been studied by \citet{badanidiyuru2013bandits} and
\citet{combes2015bandits}; we extend their framework to account for the Lipschitz
structure of PRM scores across tree depth.

\paragraph{Regret in structured environments.}
\citet{bubeck2011x} and \citet{kleinberg2008multi} study bandits with metric structure
on the action space (the $X$-armed bandit and Lipschitz bandit settings). Our
Assumption~\ref{ass:lipschitz} places DVA-MCTS in this class, allowing us to exploit
score smoothness across tree depth to reduce verifier calls while maintaining
near-optimal regret.
